\section{What it is for}
\label{sec:short:whatfor}

Everything so far has been description. This section is the argument: what a network of this shape
is good for, and why that is about to matter more than it did.

\subsection{The bet: the next tenant is software}

Open any cloud console today and you will find an interface built for a person. Pick a region. Choose
an instance size. Enter a card. Agree to terms. Every major cloud is built on the assumption that a
human is making the decision, and the whole surface --- accounts, billing relationships, credit
checks, support tiers --- follows from it.

\Flux{} is being rebuilt on the opposite assumption: that the caller is a program. Something that
provisions capacity when it needs it, pays for it, scales it, and releases it when the work is done,
with no human in the loop and no account anywhere.

The reason this is possible here and awkward everywhere else is narrow and concrete. On \Flux{},
paying for a deployment is a signed transaction. It is not a billing relationship. There is no
account to open, no card to hold, no legal entity to be. \textbf{A program can hold a key. It cannot
hold a corporate account.} That one asymmetry is the whole of the positioning claim, and everything
in this section follows from it.

This is not a prediction that people stop using clouds. It is a claim about where the growth goes.
Software that writes software, agents that run errands, pipelines that spin up a hundred short-lived
workers and vanish --- these are workloads whose natural interface is a key and a transaction, not a
console and an invoice.

\subsection{One substrate, and it already runs more than AI}

It would be easy, and wrong, to present \Flux{} as an AI network.

What actually runs on it today, by census rather than by aspiration, is a general-purpose mix:
ordinary web applications, blockchain infrastructure, content sites --- and \textbf{27 folding and
Rosetta deployments} doing real protein-folding and structural-biology work, which alone account for
roughly \SI{9}{\percent} of committed fleet CPU. Scientific simulation is not a use case \Flux{} is
hoping to attract. It is a use case that is already there, paying, at scale.

That matters more than it sounds, and the reason is economic rather than technical. A network that
serves only inference is idle whenever inference is idle. A network that serves web hosting,
scientific simulation, blockchain infrastructure and model inference --- from the same collateral,
the same placement, the same payment path --- has a demand base whose parts do not move together.
The generality is the resilience.

So the claim is not that \Flux{} becomes an AI cloud. It is that \Flux{} is a compute substrate on
which AI is one arriving workload class among several that are already running.

\subsection{Cloud and edge, converging}

\FluxCloud{} and \FluxEdge{} are to merge into a single product, offering accelerated compute through
one interface. A tenant should describe the work and let placement decide where it runs --- a
container on a commodity node, or a GPU on a datacenter machine --- without having to know which
product they are in. That convergence is planned, not shipped.

\limitnote{Convergence is at present an architectural intention with thin implementation behind it.
The only concrete artefact is a small, new capability specification, and the placement code that
would have to consume it does not read it yet: the three fields the new specification adds for this
purpose appear nowhere in the placement path. A merge announced is not a merge shipped, and this
document would rather say so than let a reader discover it later.}

\subsection{Why small models fit this fleet in particular}

The interesting question is not whether a decentralized network can run AI workloads. It is which
ones, and why.

Frontier-model inference is centralised for a specific engineering reason: the model does not fit on
one device, so it is sharded across many accelerators that must exchange activations at every layer.
That demands an interconnect --- machines physically close, wired together at bandwidths a wide-area
network cannot approach. A fleet of commodity machines scattered across the world is genuinely
unsuited to it, and no amount of decentralization enthusiasm changes that.

\textbf{But that requirement does not exist for a model that fits on a single device.} A model small
enough to sit in one machine's memory needs no interconnect at all, because there is nothing to
shard. The property that makes a dispersed fleet useless for frontier inference is simply absent for
the models most people actually need: the one that drafts the email, summarises the document,
triages the inbox, answers the question. Most work does not require a frontier model, and the work
that does not is precisely the work this shape of network can serve.

\begin{figure}[H]
\centering
\begin{tikzpicture}[font=\small,
  gpu/.style={draw=FluxSlate, fill=FluxSlate!14, rounded corners=1.5pt,
              minimum width=9mm, minimum height=11mm, inner sep=1pt,
              font=\scriptsize, text=FluxSlate},
  fnode/.style={draw=FluxBlue, fill=FluxBlue!10, rounded corners=2pt,
                minimum width=16mm, minimum height=12mm, align=center,
                font=\scriptsize, text=FluxBlue!75!black},
  lbl/.style={font=\scriptsize, text=FluxSlate},
]
%% ---- left: sharded frontier inference ----
\node[lbl, anchor=south, font=\scriptsize\bfseries] at (2.7,2.05)
  {Frontier model: sharded};
\foreach \i/\x in {0/1.0,1/2.15,2/3.3,3/4.45} \node[gpu] (g\i) at (\x,1.15) {GPU};
\foreach \i/\j in {0/1,1/2,2/3} \draw[<->, thick, FluxSlate!80] (g\i) -- (g\j);
\begin{scope}[on background layer]
  \node[draw=FluxSlate!35, dashed, rounded corners=3pt, inner sep=5pt,
        fit=(g0)(g3)] (rack) {};
\end{scope}
\node[lbl, anchor=north, text width=52mm, align=center] at (2.7,0.42)
  {activations cross the links at every layer\\
   \textbf{needs a fast interconnect}\\
   one building, or it does not run};

%% ---- divider ----
\draw[FluxRule, thick] (6.35,-0.75) -- (6.35,2.15);

%% ---- right: single-device small model ----
\node[lbl, anchor=south, font=\scriptsize\bfseries] at (10.6,2.05)
  {Small model: fits on one device};
\node[fnode] (n0) at (8.3,1.15) {node};
\node[fnode] (n1) at (10.6,1.15) {node};
\node[fnode] (n2) at (12.9,1.15) {node};
\node[lbl, anchor=north, text width=58mm, align=center] at (10.6,0.42)
  {nothing to shard, so nothing to exchange\\
   \textbf{no interconnect requirement}\\
   anywhere there is memory};
\end{tikzpicture}
\caption{Why the constraint that centralises frontier inference does not bind here. Sharding a model
across accelerators forces them into one building; a model that fits on one device imposes no such
requirement, and a dispersed fleet stops being a disadvantage.}
\label{fig:short:interconnect}
\end{figure}

There is a second, sharper reason the fit is good, and it is a property of this fleet rather than of
decentralized networks generally. Decode speed follows a model's \emph{active} parameters, but
residency --- whether the model fits at all --- follows its \emph{total} parameters. Mixture-of-experts
models separate the two: they read only a few billion parameters per token while having to hold
tens of billions in memory. Such models therefore want \textbf{RAM} more than they want memory
bandwidth. RAM is what this fleet has, in quantity, across thousands of independent machines.

On commodity dual-channel hardware a four-billion-parameter quantised model runs at roughly
\num{12} tokens per second; on an eight-channel server, roughly \num{40}. Across the network's
tenant-available memory the capacity, at full utilisation, is on the order of tens of thousands of
concurrent small-model instances.

\limitnote{Those rates are a memory-bandwidth model --- achievable bandwidth divided by bytes read
per token --- and have \emph{not} been measured on \Flux{} hardware. A benchmark on each tier's
reference machines is scheduled for Q4~2026, and until it lands these are estimates rather than
results. The capacity figures are upper bounds at full utilisation, not availability: roughly
\SI{9}{\percent} of fleet CPU is already committed to the scientific workloads above. And the GPU
path has a harder limit --- accelerators are attached whole, with no subdivision, so a tenant wanting
a fraction of a card cannot have one, and a GPU request under contention is silently dropped rather
than refused.}

\subsection{What an autonomous tenant can and cannot do}

The agent story is the most likely place for a reader to assume more than is true, so it is worth
closing on precisely.

An agent on \Flux{} can hold a key, sign a deployment specification, pay for it, watch it run and
let it expire --- with no human, no account, and no intermediary at any step. That is real today for
the payment path, and it is the direct consequence of settlement being a transaction rather than a
relationship.

What an agent cannot be given is a spending limit the network enforces. \Flux{} inherits Bitcoin's
model, in which validating a transaction is a function of that transaction and the coins it spends
--- and nothing else. A rule like \emph{spend at most five \flux{} per day} requires knowing what the
key spent yesterday, and there is nowhere for that rule to live. \textbf{An agent's budget is
whatever you fund its key with, and that is the honest bound.} Anything richer needs a co-signer,
which puts back the intermediary the design exists to remove.

That is a real constraint and the paper states it rather than working around it. It is also, in
practice, a workable one: funding a key with what you are willing to lose is a familiar discipline,
and it fails closed.
